\section{Conclusion}
\label{sec:conclusion}

In this paper, we presented a formal symbolic security analysis of the MACsec Key Agreement (MKA) protocol using the Tamarin Prover. Our model captures a two-party MKA exchange, incorporating natural-number counter abstractions for Message Numbers ($MN$) and Key Numbers ($KN$), and covers both the initial SAK distribution and a subsequent rekey round.

We formally specified and verified fourteen security properties under a standard Dolev--Yao adversary model. The analysis proves that when the long-term Connectivity Association Key (CAK) is uncompromised, MKA successfully guarantees protocol executability, mutual agreement on session parameters, SAK confidentiality, strict session ordering, and key freshness across rekeying rounds.

At the same time, our verification reveals a structural trust limitation in MKA's acceptance logic: a compromised or malicious Key Server possessing the CAK can inject an arbitrary SAK value that the Server accepts. Because non-Key Server participants rely solely on ICV verification using the derived ICK, they cannot independently verify whether an accepted SAK originated from the standard KDF path. This underlines an implicit trust assumption in IEEE 802.1X-2020 and motivates future work on cryptographic binding mechanisms for SAK derivation.

\paragraph{Future Work.} Future directions include extending the model to multi-party MKA topologies, modeling dynamic peer lists and election logic, and combining symbolic verification with computational proofs. Our complete Tamarin model is publicly available at \url{https://github.com/KaplanHalil/MKA_tamarin}.